\section{Preliminaries}\label{sec:prelim}

This section collects the tools used in the proof. Classical results are quoted; proofs are given for
the variants that are less standard.

\subsection{Kernels and operators}\label{sec:kernels}

For a kernel~$K$ we use the operator norm
\begin{equation*}
  \norm{K}\coloneqq\sup\Bigl\{\Bigl|\sum_{x,y}\overline{g(x)}K(x,y)f(y)\Bigr|\colon
  f,g\text{ finitely supported},\ \norm{f}_2=\norm{g}_2=1\Bigr\}\in[0,\infty]\,.
\end{equation*}
If it is finite, $K$ extends to a bounded operator on $\ell^2(V)$ whose matrix is~$K$; conversely, a
bounded operator $S$ has the matrix $S(x,y)=\langle\1_x,S\1_y\rangle$, where $\1_x$ is the indicator
function of~$\{x\}$. We identify the two.
For a nonnegative kernel $L$ and vertex sets $A,B$ we write
\begin{equation*}
  L(A,B)\coloneqq\sum_{a\in A}\sum_{b\in B}L(a,b)\,;
\end{equation*}
for instance, $\Adj(A,B)$ is the number of edges with one endpoint in $A$ and the other in~$B$.

\begin{lemma}[Kernel bounds]\label{lem:kernels}
Let $K$ and $L$ be kernels with $L\ge0$.
\begin{enumerate}[label=\textup{(\alph*)}]
  \item\label{it:domination} \emph{Domination.} If $\abs{K}\le L$ entrywise, then $\norm{K}\le\norm{L}$.
  \item\label{it:schur} \emph{Schur test.} If all row sums and all column sums of $L$ are at most~$b$,
  then $\norm{L}\le b$.
  \item\label{it:monotone} \emph{Monotone limits.} For $0<p_*\le1$ and all $x,y$,
  $T_p(x,y)\uparrow T_{p_*}(x,y)$ as $p\uparrow p_*$, and
  $\norm{T_{p_*}}=\sup_{p<p_*}\norm{T_p}\in[0,\infty]$. Moreover $T_p\le T_{p'}$ entrywise for
  $p\le p'$; in particular $T_p$ is bounded for every $p<\ptwo$.
\end{enumerate}
\end{lemma}

\begin{proof}
For~\ref{it:domination}, $\abs{\sum\overline{g(x)}K(x,y)f(y)}\le\sum\abs{g(x)}L(x,y)\abs{f(y)}
\le\norm{L}\norm{f}_2\norm{g}_2$. For~\ref{it:schur}, the Cauchy--Schwarz inequality with the weights
$L(x,y)$ gives
$\sum L(x,y)\abs{g(x)}\abs{f(y)}\le(\sum L(x,y)\abs{g(x)}^2)^{1/2}(\sum L(x,y)\abs{f(y)}^2)^{1/2}
\le b\norm{f}_2\norm{g}_2$.

For~\ref{it:monotone}, monotonicity in $p$ follows from the standard monotone coupling. The event
$\{x\conn y\}$ is the increasing union of the events $\{x\connin{B(x,n)}y\}$, $n\in\N$, whose
probabilities are polynomials in~$p$; hence $\lim_{p\uparrow p_*}T_p(x,y)=\sup_p\sup_n
\PP_p(x\connin{B(x,n)}y)=\sup_n\PP_{p_*}(x\connin{B(x,n)}y)=T_{p_*}(x,y)$. If
$\norm{T_p}\le M$ for all $p<p_*$, then for finitely supported $f,g$ the finite sum
$\sum\overline{g(x)}T_{p_*}(x,y)f(y)$ is the limit of the corresponding sums for $T_p$, so
$\norm{T_{p_*}}\le M$. The reverse inequality and the last assertion follow from monotonicity
and~\ref{it:domination}.
\end{proof}

\subsection{Unimodularity, mass transport and the trace}\label{sec:mtp}

Equip $\Aut(G)$ with the topology of pointwise convergence. It is a locally compact group whose vertex
stabilisers $\operatorname{Stab}(x)\coloneqq\{g\in\Aut(G)\colon gx=x\}$ are compact and open. The group
is \emph{unimodular} if its left Haar measure is also right-invariant. For a transitive graph this is
equivalent to $\abs{\operatorname{Stab}(x)y}=\abs{\operatorname{Stab}(y)x}$ for all $x,y\in V$, where
$\operatorname{Stab}(x)y$ is the orbit of $y$ under $\operatorname{Stab}(x)$, and it is equivalent to the mass-transport
principle below; see~\cite[Section~3]{BLPS99b} and~\cite[Chapter~8]{LP16}. For instance, the
automorphism group of every Cayley graph is unimodular~\cite[Chapter~8]{LP16}. A function
$F\colon V\times V\to[0,\infty]$ is \emph{diagonally invariant} if $F(gx,gy)=F(x,y)$ for all
$g\in\Aut(G)$ and $x,y\in V$.

\begin{lemma}[Mass-transport principle, \cite{BLPS99b}, {\cite[Chapter~8]{LP16}}]\label{lem:mtp}
Assume \cref{ass:G,ass:U}. For every diagonally invariant $F\colon V\times V\to[0,\infty]$,
\begin{equation*}
  \sum_{y\in V}F(o,y)=\sum_{x\in V}F(x,o)\,.
\end{equation*}
\end{lemma}

We think of $F(x,y)$ as mass sent from $x$ to~$y$: the expected mass sent out by the root equals the
expected mass it receives. The identity extends to diagonally invariant real or complex $F$ with
$\sum_y\abs{F(o,y)}<\infty$, by splitting into positive and negative parts. Typically
$F(x,y)=\EE_p[m(x,y;\omega)]$ for a random mass $m$ that is \emph{equivariant}, that is,
$m(gx,gy;g\omega)=m(x,y;\omega)$, where $g\omega$ is the image of the configuration under~$g$; such
$F$ are diagonally invariant because $\PP_p$ is invariant. Two consequences are used repeatedly.
First, every nonnegative invariant kernel has equal row and column sums; in particular an invariant
stochastic kernel is doubly stochastic and has norm at most one by \cref{lem:kernels}. Second, the
diagonal of an invariant operator defines a trace.

A bounded operator $S$ on $\ell^2(V)$ is \emph{invariant} if its matrix is an invariant kernel,
equivalently if $S$ commutes with the unitary operators $U_gf\coloneqq f\circ g^{-1}$, $g\in\Aut(G)$.
For such $S$ we set
\begin{equation*}
  \tr S\coloneqq S(o,o)=\langle\1_o,S\1_o\rangle\,.
\end{equation*}
Since $\1_o$ is a unit vector, $\abs{\tr S}\le\norm S$. For example, $\tr(T_p^3)=\nabla_p$ whenever
$T_p$ is bounded.

\begin{lemma}[The trace]\label{lem:trace}
Assume \cref{ass:G,ass:U}, and let $A$ and $B$ be bounded invariant operators.
\begin{enumerate}[label=\textup{(\alph*)}]
  \item\label{it:trace-cyclic} $\tr(AB)=\tr(BA)$.
  \item\label{it:trace-positive} If $A$ is positive semidefinite and $B$ is self-adjoint, then
  $\abs{\tr(BA)}\le\norm{B}\tr A$.
  \item\label{it:trace-diagonal} If $A$ is positive semidefinite, then $A(x,x)=\tr A$ and
  $\abs{A(x,y)}\le\tr A$ for all $x,y\in V$.
\end{enumerate}
\end{lemma}

\begin{proof}
For~\ref{it:trace-cyclic}, the series defining $\tr(AB)$ converges absolutely, since
\begin{equation*}
  \sum_y\abs{A(o,y)B(y,o)}\le\norm{A^*\1_o}_2\,\norm{B\1_o}_2\le\norm{A}\,\norm{B}\,.
\end{equation*}
The function $F(x,y)\coloneqq A(x,y)B(y,x)$ is diagonally invariant, and the mass-transport principle
gives $\tr(AB)=\sum_yA(o,y)B(y,o)=\sum_xA(x,o)B(o,x)=\tr(BA)$.

For~\ref{it:trace-positive}, the positive square root $A^{1/2}$ commutes with every $U_g$, because $A$
does and the functional calculus preserves commutation; so $A^{1/2}$ is invariant. By~\ref{it:trace-cyclic},
$\tr(BA)=\tr(A^{1/2}BA^{1/2})$. Since $-\norm{B}A\le A^{1/2}BA^{1/2}\le\norm{B}A$ in the order of
quadratic forms, evaluating at $\1_o$ gives the claim.

For~\ref{it:trace-diagonal}, $A(x,x)=A(o,o)$ by invariance and transitivity, and
$\abs{A(x,y)}^2\le A(x,x)A(y,y)$ by the Cauchy--Schwarz inequality for the positive semidefinite form
$(f,g)\mapsto\langle f,Ag\rangle$.
\end{proof}

\subsection{The lazy random walk}\label{sec:walk}

Recall from \cref{sec:mechanism} the lazy simple random walk $P=(I+\Adj/d)/2$ and
$\rho=\norm{P}$. The following is a version of the discrete Cheeger inequality; see~\cite{Dod84} for its
background.

\begin{lemma}[Spectral gap]\label{lem:walk}
Assume \cref{ass:G}. The kernel $P$ is symmetric, invariant and positive semidefinite, and
\begin{equation*}
  \rho=\norm{P}\le1-\frac{\iota^2}{4d^2}<1\,.
\end{equation*}
\end{lemma}

\begin{proof}
Symmetry and invariance are clear. By the Schur test $\norm{\Adj/d}\le1$, so $P=(I+\Adj/d)/2$ is
positive semidefinite. For finitely supported $f$,
\begin{equation*}
  \langle f,(I-P)f\rangle=\frac1{2d}\sum_{\{x,y\}\in E}\abs{f(x)-f(y)}^2\,.
\end{equation*}
Let $f\ge0$. By the co-area formula and the isoperimetric inequality,
$\sum_{\{x,y\}\in E}\abs{f(x)^2-f(y)^2}=\int_0^\infty\abs{\partial_E\{f^2>s\}}\dd s\ge\iota\norm{f}_2^2$.
The Cauchy--Schwarz inequality and $\sum_{\{x,y\}\in E}(f(x)+f(y))^2\le2d\norm{f}_2^2$ bound the
left-hand side by $(2d)^{1/2}\norm{f}_2(\sum_{\{x,y\}\in E}(f(x)-f(y))^2)^{1/2}$. Hence
$\langle f,(I-P)f\rangle\ge\iota^2\norm{f}_2^2/(4d^2)$, that is,
$\langle f,Pf\rangle\le(1-\iota^2/(4d^2))\norm{f}_2^2$. Since $P$ has nonnegative entries,
$\abs{\langle f,Pf\rangle}\le\langle\abs{f},P\abs{f}\rangle$ for general~$f$, and for the self-adjoint
operator $P$ this bound on the quadratic form bounds the norm.
\end{proof}

\subsection{Correlation inequalities}\label{sec:correlation}

An event or a function of the configuration is \emph{increasing} if it is preserved, respectively does
not decrease, when edges are opened. Two increasing events $\mathcal A,\mathcal B$ occur
\emph{disjointly}, written $\mathcal A\circ\mathcal B$, if there are disjoint finite sets of open edges
that witness $\mathcal A$ and $\mathcal B$ respectively.

\begin{lemma}[Harris, van den Berg--Kesten, tree-graph bound]\label{lem:correlation}
Let $p\in[0,1]$.
\begin{enumerate}[label=\textup{(\alph*)}]
  \item\label{it:harris} Increasing nonnegative functions of the configuration are positively
  correlated under $\PP_p$.
  \item\label{it:bk} For increasing events $\mathcal A_1,\dots,\mathcal A_m$ whose occurrence is
  witnessed by finite sets of open edges,
  $\PP_p(\mathcal A_1\circ\dots\circ\mathcal A_m)\le\prod_i\PP_p(\mathcal A_i)$.
  \item\label{it:tree} For all $x,y,z\in V$ and every domain $\Lambda$,
  \begin{equation*}
    \PP_p(x\connin{\Lambda}y,\ x\connin{\Lambda}z)\le\sum_{w\in V}T_p(x,w)\,T_p(w,y)\,T_p(w,z)\,.
  \end{equation*}
  In particular $\EE_p\abs{C_x}^2\le\chi(p)^3$.
\end{enumerate}
\end{lemma}

\begin{proof}
Parts~\ref{it:harris} and~\ref{it:bk} are the Harris inequality~\cite{Har60} and the
van den Berg--Kesten inequality~\cite{vdBK85}; see~\cite[Sections~2.2 and~2.3]{Gri99}. The passage from
finitely many edges to events with finite witnesses on~$G$ is by monotone approximation. For~\ref{it:tree},
on the event on the left choose self-avoiding open paths in $\Lambda$ from $x$ to $y$ and from $x$ to~$z$.
Their union contains an open tree spanning $x,y,z$, which has a vertex $w$ joined to $x$, $y$ and $z$ by
three edge-disjoint open paths; paths of length zero are allowed. A union bound over $w$ and
part~\ref{it:bk} give the claim, using $T_p^\Lambda\le T_p$. Summing over $y,z$ gives the bound on the
second moment.
\end{proof}

\subsection{The subcritical regime}\label{sec:subcritical}

\begin{lemma}[Subcritical susceptibility]\label{lem:subcritical}
Assume \cref{ass:G}. For every $p<p_c$ the susceptibility $\chi(p)$ is finite. The kernel $T_p$ is then
a bounded, invariant, symmetric and positive semidefinite operator with $\norm{T_p}\le\chi(p)$.
\end{lemma}

\begin{proof}
Finiteness of $\chi(p)$ below $p_c$ is the sharpness of the phase transition. It was proved on the
hypercubic lattice by Aizenman and Barsky~\cite{AB87} and Menshikov~\cite{Men86}, and on quasi-transitive
graphs by Antunovi\'c and Veseli\'c~\cite[Theorem~2]{AV08}; on every locally finite, transitive, infinite
graph it is~\cite[Theorem~1.1]{DCT16}, which also gives a short proof. The norm bound is the Schur test. Since
$T_p(x,y)=\EE_p\sum_C\1\{x,y\in C\}$, where the sum runs over the clusters, every finitely supported $f$
satisfies $\sum_{x,y}\overline{f(x)}T_p(x,y)f(y)=\EE_p\sum_C\abs{\sum_{x\in C}f(x)}^2\ge0$.
\end{proof}

The next lemma is Russo's formula in the form we need: in operator norm and for the infinite graph. Its
proof is given in \cref{app:russo}.

\begin{lemma}[Russo's formula]\label{lem:russo}
Assume \cref{ass:G,ass:U}. The map $p\mapsto T_p$ is locally absolutely continuous on $(0,p_c)$ in
operator norm, and $\chi$ is locally absolutely continuous on $(0,p_c)$. For almost every $p\in(0,p_c)$,
\begin{align}
  (1-p)\,T_p'(x,y)&=\EE_p\bigl[\Adj(C_x,C_y)\,\1\{C_x\neq C_y\}\bigr]\qquad\text{for all }x,y\in V\,,
  \label{eq:russo-kernel}\\
  (1-p)\,\chi'(p)&=\sum_{b\sim o}\EE_p\bigl[\abs{C_o}\,\abs{C_b};\,C_o\neq C_b\bigr]
  =\chi(p)^2\,\Delta(p)\,,\label{eq:russo-chi}
\end{align}
where $T_p'$ is the derivative in operator norm and $\Delta(p)$ is defined by the last identity.
\end{lemma}

The quantity $\Delta(p)$ is the normalised expected product of the sizes of two adjacent distinct
clusters; it already appeared in \cref{sec:part1}. In~\eqref{eq:russo-kernel}, $\Adj(C_x,C_y)$ counts
the closed edges between the two clusters, each of which is pivotal for $x\conn y$.

\begin{lemma}[Sprinkling, {\cite[Lemma~2.4]{Hut19}}]\label{lem:sprinkling}
For $p,a\in[0,1]$,
\begin{equation}\label{eq:sprinkling}
  T_{p+(1-p)a}\le\sum_{j\ge0}a^j\,T_p(\Adj\,T_p)^j\qquad\text{entrywise.}
\end{equation}
Consequently, if $ad\norm{T_p}<1$ then $\norm{T_{p+(1-p)a}}\le\norm{T_p}/(1-ad\norm{T_p})$, and if
$ad\chi(p)<1$ then $\chi(p+(1-p)a)\le\chi(p)/(1-ad\chi(p))$.
\end{lemma}

\begin{proof}
Let $\omega_1$ and $\omega_2$ be independent configurations with laws $\PP_p$ and $\PP_a$. Their union
$\omega_1\vee\omega_2$ has law $\PP_{p+(1-p)a}$. If $x$ and $y$ are joined by a self-avoiding path open
in $\omega_1\vee\omega_2$, let $e_1,\dots,e_j$ be the edges of the path, in order, that are closed
in~$\omega_1$; they are open in $\omega_2$ and pairwise distinct. Orient $e_i=(x_i,y_i)$ along the path.
The pieces of the path between them are edge-disjoint and open in~$\omega_1$, joining $x$ to $x_1$,
$y_1$ to $x_2$, \dots, $y_j$ to~$y$. By a union bound over $j$ and the oriented edges, independence of
$\omega_2$ and the van den Berg--Kesten inequality for~$\omega_1$,
\begin{equation*}
  T_{p+(1-p)a}(x,y)\le\sum_{j\ge0}a^j\sum_{(x_i,y_i)_{i\le j}}T_p(x,x_1)\,T_p(y_1,x_2)\cdots T_p(y_j,y)\,,
\end{equation*}
which is~\eqref{eq:sprinkling}. Since $\norm{\Adj}\le d$ and the row sums of $\Adj$ equal~$d$, the
norm and the row sums of the right-hand side are at most $\sum_ja^jd^j\norm{T_p}^{j+1}$ and
$\sum_ja^jd^j\chi(p)^{j+1}$, respectively. \Cref{lem:kernels}\ref{it:domination} concludes.
\end{proof}

\subsection{Criticality and two reductions}\label{sec:criticality}

\begin{lemma}[Critical finiteness]\label{lem:pc}
Assume \cref{ass:G}. Then $0<p_c<1$, and $\PP_{p_c}$-almost surely every cluster is finite. In
particular, for every domain $\Lambda$, all clusters $C^\Lambda_x$ are finite $\PP_{p_c}$-almost surely.
\end{lemma}

\begin{proof}
There are at most $d^n$ self-avoiding paths of length $n$ from~$o$, so $\theta(p)=0$ for $p<1/d$. On the
other hand, a connected vertex set $A\ni o$ with $\abs{A}=n$ can be encoded by a closed walk of length
$2n-2$ from $o$ around a spanning tree of~$A$, so there are at most $d^{2n}$ such sets. If $C_o=A$, then
all $\abs{\partial_E A}\ge\iota n$ edges of $\partial_EA$ are closed. Hence
$\PP_p(\abs{C_o}<\infty)\le\sum_{n\ge1}d^{2n}(1-p)^{\iota n}<1$ for $p$ close to one, and $p_c<1$.

Every ball satisfies $\abs{B(o,r+1)}\ge(1+\iota/d)\abs{B(o,r)}$, since each vertex of the outer vertex
boundary of $B(o,r)$ is incident to at most $d$ edges of $\partial_EB(o,r)$. Thus $G$ has exponential
growth, and by a theorem of Hutchcroft~\cite[Theorem~1.2]{Hut16} there are almost surely no infinite
clusters at~$p_c$. Clusters in a domain are contained in clusters of~$G$.
\end{proof}

\begin{lemma}[The operator threshold precedes uniqueness]\label{lem:p22-pu}
Assume \cref{ass:G}. If $\PP_p(N_\infty=1)=1$, then $T_p$ is unbounded. Consequently $\ptwo\le p_u$.
\end{lemma}

\begin{proof}
If $\PP_p(N_\infty=1)=1$, then some vertex lies in an infinite cluster with positive probability, so
$\theta(p)>0$ by transitivity. By uniqueness and the Harris inequality,
$T_p(x,y)\ge\PP_p(\abs{C_x}=\abs{C_y}=\infty)\ge\theta(p)^2$ for all $x,y$. For finite $F\subseteq V$
this gives $\langle\1_F,T_p\1_F\rangle\ge\theta(p)^2\abs{F}^2=\theta(p)^2\abs{F}\,\norm{\1_F}_2^2$, so
$T_p$ is unbounded. If we had $\ptwo>p_u$, the definition of $p_u$ as an infimum would provide
$p<\ptwo$ with $\PP_p(N_\infty=1)=1$, and $T_p$ would be both bounded (\cref{lem:kernels}) and
unbounded.
\end{proof}

\begin{proposition}[The nonunimodular case, {\cite[Theorem~2.9]{Hut19}}]\label{prop:nonunimodular}
Assume \cref{ass:G}. If $\Aut(G)$ is not unimodular, then $p_c<\ptwo$.
\end{proposition}

\begin{proof}
Hutchcroft's theorem states that $p_c<p_{q\to q}$ for every $q\in(1,\infty)$ on every connected, locally
finite graph whose automorphism group has a quasi-transitive nonunimodular subgroup. Under our
assumption, $\Aut(G)$ itself is such a subgroup, and $q=2$ gives the claim.
\end{proof}

The last result of this section turns boundedness of one connection operator above $p_c$ into a
simultaneous nonuniqueness statement. It follows the argument of~\cite{Ope26}.

\begin{proposition}[Simultaneous nonuniqueness]\label{prop:interval}
Assume \cref{ass:G}, and let $p_c<p_1<p_2<1$ be such that $T_{p_2}$ is bounded. In the standard
monotone coupling, almost surely for every $p\in[p_1,p_2]$ there are infinitely many infinite clusters
at parameter~$p$.
\end{proposition}

\begin{proof}
\emph{Stage 1: ergodicity.} Every $\Aut(G)$-invariant event $\mathcal E$ of the labels $(U_e)_{e\in E}$
has probability zero or one. Indeed, for $r\in\N$ choose an event $\mathcal A_r$ depending only on the
labels in the ball $B(o,r)$ such that $\PP(\mathcal E\triangle\mathcal A_r)\to0$ as $r\to\infty$, and an
automorphism $g_r$ with $\dist(o,g_ro)>2r+2$, which exists by transitivity. Then $\mathcal A_r$ and
$g_r\mathcal A_r$ are independent, while $g_r\mathcal E=\mathcal E$, so
$\abs{\PP(\mathcal E)-\PP(\mathcal E)^2}\le4\PP(\mathcal E\triangle\mathcal A_r)\to0$.

\emph{Stage 2: the count.} Let $N_{12}\in\{0,1,2,\dots\}\cup\{\infty\}$ be the number of
$p_2$-clusters that contain an infinite $p_1$-cluster. It is invariant, hence almost surely constant.
Since $\theta(p_1)>0$, it is positive.

It is not one. Otherwise any two vertices in infinite $p_1$-clusters would lie in the same
$p_2$-cluster, and the Harris inequality would give $T_{p_2}(x,y)\ge\theta(p_1)^2$ for all $x,y$; as
in the proof of \cref{lem:p22-pu}, $T_{p_2}$ would be unbounded.

It is not an integer $n\ge2$. Otherwise there would be vertices $x,y$ such that, with positive
probability, $x$ and $y$ lie in infinite $p_1$-clusters contained in distinct $p_2$-clusters; call
this event~$\mathcal E_{xy}$. Fix a path $\pi$ from $x$ to $y$, and replace the labels of the edges of
$\pi$ by independent labels uniform on $[0,p_1]$. The law of the modified labels is absolutely
continuous with respect to the original law, since finitely many coordinates were replaced by
variables with bounded densities. The modification only opens edges, at both parameters, and only
finitely many. Hence clusters at either parameter can only merge: every cluster after the modification
is the union of finitely many clusters before it. In particular every infinite $p_1$-cluster after the
modification contains an infinite $p_1$-cluster from before. Consequently every counted $p_2$-cluster
before the modification lies in a counted $p_2$-cluster after it, and every counted $p_2$-cluster after
the modification contains a counted one from before. On $\mathcal E_{xy}$ the path $\pi$ joins two
distinct counted $p_2$-clusters, which therefore lie in the same counted cluster after the modification.
So the modified labels have $N_{12}\le n-1$ with positive probability, which contradicts absolute
continuity and $N_{12}=n$ almost surely.

\emph{Stage 3: conclusion.} Hence $N_{12}=\infty$ almost surely. On this event choose infinite
$p_1$-clusters in infinitely many distinct counted $p_2$-clusters. For every $p\in[p_1,p_2]$ they lie in
infinite $p$-clusters, which are distinct because they lie in distinct $p_2$-clusters.
\end{proof}
